\selectlanguage{ngerman}
\Language{de}{Deutsch}{Benutzerhandbuch}
\firsttopic{Öffnen, schlafen, zurückkehren}
Inkline führt auf dem reMarkable 2 eine lokale Shell aus. Verwenden Sie ein
Type Folio oder eine USB-Tastatur, arbeiten Sie direkt auf dem Tablet oder
verbinden Sie sich über SSH oder Goblin Mosh. Text und Bilder erscheinen in Graustufen.

\begin{notice}
\key{Strg+Opt+Alt+T} öffnet Inkline direkt auf dem Tablet.\par
\key{Ein/Aus-Taste:} drücken und loslassen, um den Ruhezustand zu starten;
erneut drücken zum Aufwachen. Shells und Editorpuffer bleiben im RAM.
\end{notice}
Der Ruhezustand pausiert das Tablet, ohne Inkline zu beenden oder Einstellungen
auf Flash zu schreiben. Netzwerkverbindungen müssen sich eventuell neu verbinden.
Der Tastendruck zum Aufwachen löst keinen erneuten Ruhezustand aus. Nach einem
Update Inkline einmal schließen und öffnen, damit die neue Tastenbehandlung gilt.

Zurück zu den Notizbüchern: \key{Quit} antippen oder \key{Strg+Umschalt+Q}
drücken und bestätigen. \code{exit} schließt ein Terminal; das letzte führt
zurück zu den Notizbüchern. Unten stehen Escape, Verlauf auf/ab, Quit und Settings.
\code{clear} entfernt die kurze Starthilfe und das kleine Goblin-Logo.

Belegt eine per Kürzel gestartete App den Bildschirm, führt \key{Strg+Opt+Alt+T}
zu Inkline zurück. Im Notfall \key{Strg+Opt+Alt+Backspace zwei Sekunden halten}.
Die Wiederherstellung schließt alle Terminals; ungespeicherte Arbeit darin geht
verloren. Nach einem Neustart erscheint weiterhin die normale Notizbuchoberfläche.

\ProjectLinks

\topic{Tasten und Terminalplätze}
Das Type Folio hat eine eigene \key{Opt}-Taste zwischen Ctrl und Alt sowie
eine \key{rechte Alt/Opt}-Taste. Sie haben unterschiedliche Aufgaben.
Ctrl heißt hier Strg, Shift heißt Umschalt.

\begin{keytable}
\key{Tasten} & \key{Aktion} \\ \midrule
Separate Opt + 1--0 & F1--F10 \\
Rechte Alt/Opt + 1--9 & Terminalplatz 1--9 wählen \\
Rechte Alt/Opt + Links / Rechts & Voriger / nächster Platz \\
Rechte Alt/Opt + Tab & Escape \\
Rechte Alt/Opt + Auf / Ab & Bild auf / Bild ab \\
Rechte Alt/Opt + Backspace & Schließen aller Terminals bestätigen \\
Beliebige Alt + Leertaste & Unicode-Tastatur \\
Rechte Alt/Opt + C / V & Auswahl kopieren / einfügen \\
Strg+Umschalt+B & Untere Leiste aus- / einblenden \\
\end{keytable}

Auf dem \key{US-Type-Folio} erzeugt rechte Alt/Opt+\key{0} ein \code{+}.
Rechte Alt/Opt mit der \key{Minustaste direkt links von Backspace} erzeugt
\code{=}. Bei ausgeschalteter Eingabemethode liefert Umschalt+6 ein wörtliches
Zirkumflex: \code{c^2} bleibt \code{c^2}. Es gibt keine Zoom-Tastenkürzel.

USB-Tastaturen behalten ihre üblichen Funktionstasten. Das Kürzel für die
separate Opt-Taste verwendet Meta, wenn die USB-Tastatur diese Taste bereitstellt.

Es gibt höchstens \key{neun unabhängige Plätze}. Ein leerer Platz öffnet eine
Shell. Die Anzeige zählt offene Terminals: Sind nur 1 und 9 offen, zeigt Platz 9
\key{Terminal 2 / 2 · slot 9}. Schließen gibt den Platz frei; die Nummern bleiben stabil.

Jedes Terminal hat Shell, Verzeichnis, laufende Eingabe, Bilder und 500
Verlaufszeilen für sich. Anzeige- und Eingabemethode gelten für alle Plätze.

\topic{Einstellungen und Darstellung}
Tippen Sie auf \key{Settings}, die fünfte Schaltfläche unten. Bei ausgeblendeter
Leiste: \key{Alt+Leertaste}, dann \key{F2} oder Settings auf der Unicode-Tastatur.
Eine gefüllte Auswahl ist aktiv; eine gestrichelte Umrandung zeigt den Tastaturfokus.
Tab bewegt den Fokus, Pfeile und Enter ändern die Auswahl.

\key{Caps Lock:} Control (Standard) oder Feststelltaste wählen.
\key{Schriftgröße:} mit zwei Fingern zoomen oder Minus/Plus in Settings nutzen.
Langsame Bewegungen ändern jeweils einen Pixel, von \key{6 bis 48 Pixeln}.

\key{Text darkness:} 50\% ist die ursprüngliche Darstellung. Höhere Werte
verdunkeln die Buchstabenränder. \key{Minimum contrast} vergrößert bei ähnlichen
Farben den Helligkeitsabstand von Vorder- und Hintergrund; Standard ist 35\%.

\begin{choices}
\key{E-Paper-Modus} & \key{Abwägung} \\ \midrule
Fast & Kaum Bündelungsverzögerung; Animationswellenform; mehr Geisterbilder \\
Balanced & UI-Wellenform; Ausgabe über 32 ms bündeln \\
Crisp (Standard) & Sauberere Graustufen; Inhaltswellenform; langsamer \\
Mono & Schnelles Schwarzweiß; keine Bildgraustufen \\
Saver & Über 120 ms bündeln; weniger Aktualisierungen bei Ausgabeschüben \\
\end{choices}

Nur geänderte Textzeilen werden neu gezeichnet. Das Panel benötigt dennoch Zeit.
Die Akkuersparnis von Saver wurde nicht gemessen. Updates behalten einen
ausdrücklich gespeicherten Modus bei.

\begin{notice}
Alle Einstellungen bleiben im \key{RAM} und werden beim normalen Beenden
gemeinsam gespeichert. Fingerheben, Schließen von Settings und Ruhezustand
schreiben nicht auf Flash. Unveränderte Sitzungen schreiben nichts. Absturz,
erzwungenes Beenden oder Stromausfall können ungespeicherte Änderungen verlieren.
\end{notice}

\topic{Unicode und Eingabemethoden}
\key{Alt+Leertaste} öffnet die Unicode-Tastatur. Ein Zeichen antippen, um es
einzufügen; die Tastatur bleibt offen. \key{Done} oder Escape schließt sie.

Die Kategorien umfassen Buchstaben, kombinierende Zeichen, Zahlen, Satzzeichen,
Mathematik/Währungen, Symbole, Leerraum, Formatzeichen und private Nutzung.
Vertikal ziehen oder Pfeile, Bild auf/ab, Home/End und das Mausrad verwenden.
\key{Tab} wechselt die Kategorie, Enter fügt die Auswahl ein.

Ein hexadezimaler Unicode-Wert wie \code{03bb}, gefolgt von Enter, fügt λ ein.
Backspace bearbeitet den Wert. Der gepunktete Kreis bei kombinierenden Zeichen
ist nur eine Hilfe; eingefügt wird die Marke allein. Leerraum und Formatzeichen
haben Beschriftungen. Der Katalog folgt den Unicode-Tabellen von Qt, ohne
Steuerzeichen, unbelegte Werte, Surrogate und Nichtzeichen. Manche Glyphen
benötigen zusätzliche Schriften; private Zeichen erfordern eine passende Schrift.

In der Auswahl \key{F2} oder \key{Settings}, dann \key{Input method} wählen:
\begin{choices}
Off — direct keyboard & Normale US-Tastatureingabe \\
Romaji & Japanische Hiragana aus lateinischer Umschrift \\
Pinyin & Chinesisch über Pinyin \\
Zhuyin & Chinesisch mit grundlegender Zhuyin-Belegung \\
Wubi 86 & Chinesische Formcodes \\
US-International & Häufige lateinische Akzente \\
\end{choices}

Ein Ergebnis in der Kandidatenleiste wählen. Mit \key{Off — direct keyboard}
im selben Menü kehren Sie zur normalen US-Eingabe zurück. Die Auswahl gilt
für alle Terminals und wird beim Beenden gespeichert.

\topic{Berührung, Stift und Zwischenablage}
\key{Zwei Finger auf/ab:} im Verlauf blättern. Nach unten ziehen zeigt ältere
Ausgabe, nach oben geht es zur Eingabeaufforderung. Jedes Terminal hält
zusätzlich zum sichtbaren Bildschirm 500 physische Zeilen im RAM.

\key{Zwei Finger links/rechts:} pro Geste einen Terminalplatz wechseln.
Vor dem nächsten Wechsel die Finger anheben. Spreizen oder Zusammenziehen
ändert die Schriftgröße.

\key{Ein Finger:} ziehen, um Text auszuwählen. Rechte Alt/Opt+C kopiert,
rechte Alt/Opt+V fügt ein. Alle Terminals teilen eine RAM-Zwischenablage mit
höchstens 128 KiB.

\key{Stift:} Apps mit aktivierten Mausmeldungen erhalten Druck, Bewegung und
Loslassen als Ereignisse der linken Maustaste, mit dem ausgehandelten Protokoll
einschließlich SGR. Sonst wählt Ziehen Text aus. \key{Umschalt} erzwingt lokale
Textauswahl auch in einer App, die Mausmeldungen verwendet.

\key{OSC 52} erlaubt Terminalprogrammen, auch über SSH, diese Zwischenablage
zu lesen oder zu ersetzen. Sie ist von der Notizbuch-Zwischenablage getrennt.
Rechte Alt/Opt+X kopiert und erklärt, dass Löschen im Editor erfolgen muss:
OSC 52 kann keinen Text im Editor entfernen.

\key{Grafik:} Kitty-Übertragungen inline und über gemeinsamen Speicher bleiben
im RAM. Datei- und temporäre Dateiübertragung sind abgeschaltet. Sixel funktioniert
bei ausdrücklicher Anforderung, wird aber nicht automatisch angekündigt.
Kitty-Platzhalter, Animationen und einige Sixel-Modi sind noch unvollständig.

\code{clear} und vollständiges Bildschirmlöschen entfernen sichtbare Bilder
und Text. Bei Speicherdruck werden Bilder außerhalb des Bildschirms freigegeben.
Es gibt keinen Bildcache auf dem Datenträger; große Bilder benötigen dennoch RAM.

\topic{Programme mit Tastenkürzeln starten}
Kürzel in einer Inkline-Shell registrieren. Argumente werden wörtlich übergeben;
für Shellsyntax eine Shell ausdrücklich aufrufen. Kürzelprogramme lesen
\code{~/.bashrc} über Bash.

Unter Settings → Command shortcuts (Seite 2) lässt sich einer Buchstabentaste ein Befehl zuweisen. Wähle Terminal oder Native e-paper app und dann Save. Remove verlangt eine Bestätigung; Reset draft verwirft Änderungen. T ist gesperrt. Nur Save und bestätigtes Entfernen schreiben auf den Speicher. Das Kommandozeilenwerkzeug verwaltet dieselben Belegungen.

Alle globalen Tastenkürzel verwenden jetzt Strg+Opt+Alt, auch T und die Wiederherstellung durch Halten von Backspace. Verwende die separate Opt-Taste des Folio, auf USB-Tastaturen Windows/Command (Super). Strg+Alt bleibt für Emacs verfügbar. Beende und öffne Inkline nach dem Update erneut; die Installation erhält offene Terminals.

\begin{shell}
~/inkline shortcut register e emacs
~/inkline shortcut register p python3
~/inkline shortcut list
~/inkline shortcut deregister p
\end{shell}

\key{Strg+Opt+Alt+E} öffnet Emacs auf einem freien Platz.
\code{~/inkline run PROGRAM ARG...} startet ohne registriertes Kürzel.
Erneutes Registrieren ersetzt die Belegung; Entfernen lässt laufende Programme
unberührt. Tasten entsprechen physischen US-Positionen. Satzzeichen in
Anführungszeichen setzen und für eigene Programme vollständige Pfade verwenden.
Die Programme müssen installiert sein.

Für eine native Qt/E-Paper-App, die selbst auf den Bildschirm zeichnet:
\begin{shell}
~/inkline shortcut register a --epaper /home/root/my-app
\end{shell}
Es läuft jeweils eine native App; Inkline und seine Shells pausieren währenddessen.
Der Starter setzt Qt-Optionen und RAM-Caches, verwirft Protokolle und setzt
Inkline nach dem Ende fort. Apps sollen Qt-Eingaben ohne exklusiven Tastaturzugriff nutzen.

\begin{notice}
\key{Strg+Opt+Alt+T ist fest reserviert.} Registrieren und Entfernen sind gesperrt.
Es beendet die native App und kehrt zu bestehenden Terminals zurück.
Die physische Strg-Taste verwenden; Caps-Lock-Umbelegung gilt nur innerhalb von Inkline.
\end{notice}

\key{Notfall:} Strg+Opt+Alt+Backspace zwei Sekunden halten, um die App samt
Kindprozessen zu stoppen und Inkline neu zu starten. Alle Terminals schließen;
ungespeicherte Arbeit geht verloren. Root-Apps können die Rettung durch
exklusiven Eingabezugriff, Dienständerungen oder Ressourcenmangel verhindern.
Ein Geräteneustart bleibt der letzte Ausweg.

\topic{Installieren, aktualisieren, Handbuch}
Das Paket ist für \key{reMarkable 2 mit Firmware 3.27} und die vorhandenen
Qt-6.8-E-Paper-Bibliotheken vorgesehen. Downloads und Prüfsummenverfahren:

\InstallLink

Die Installation prüft Hardware, Firmware, Schriften, Qt-Start, temporäre
RAM-Verzeichnisse und deaktivierten Swap. Versionen liegen unter
\path{/home/root/.local/share/inkline}; der Tastaturstarter läuft ab Systemstart.
Updates erhalten offene Terminals. \key{Inkline schließen und erneut öffnen,
um die neue Version zu verwenden.} Frühere Versionen bleiben für eine Rückkehr
erhalten und belegen zusätzlichen Speicherplatz.

Dieses eine PDF enthält alle neun Sprachen und liegt dem Paket bei.
Die Installation versucht den Import nach \key{My files} über die aktivierte
USB-Weboberfläche, wenn die Notizbücher laufen. Ein erfolgreicher Import wird
vermerkt, damit dasselbe Handbuch nicht doppelt importiert wird.

Falls der Import nicht möglich ist: USB anschließen, Inkline beenden und
\key{USB web interface} in den Speichereinstellungen der Notizbücher aktivieren.
Dann per SSH ausführen:
\begin{shell}
~/inkline manual
\end{shell}
Alternativ das PDF herunterladen und über die USB-Browseroberfläche oder die
reMarkable-App für Computer/Mobilgerät importieren. Der Import verändert die
Notizbuchdatenbank nie direkt und unterbricht keine Terminals. Das importierte
Handbuch bleibt nach der Deinstallation in My files.

Per SSH zu den Notizbüchern zurückkehren oder Inkline entfernen:
\begin{shell}
~/inkline stop
~/inkline uninstall
\end{shell}
Optionale Software einschließlich der empfohlenen LaTeX-Installation wird
auf der Website beschrieben.

\topic{Funktionsweise, Speicher und RAM}
Qt übergibt Tasten an die gemeinsame Zuordnung für Kürzel und Caps Lock.
libghostty-vt codiert Terminaltasten. Jeder Platz verbindet sein Programm
mit einem eigenen Pseudoterminal (PTY). Shells erhalten
\code{HOME=/home/root}; interaktives Bash liest \code{~/.bashrc}.

Ausgabe durchläuft den Terminalparser und Sixel-Decoder. libghostty verwaltet
Zellen, Cursor, alternative Bildschirme, Verlauf und Kitty-Bilder. Der Renderer
hält ein Graustufenbild und zeichnet geänderte Zeilen neu; das Qt-E-Paper-Backend
aktualisiert das Panel. Ein eigener evdev-Daemon verarbeitet globale Kürzel
und die Ein/Aus-Taste; systemd übernimmt Schlafen und Aufwachen.

Pro Terminal sind Kernzuweisungen auf 64 MiB begrenzt, Kitty-Bilder auf
16 MiB je Bildschirm und aufbewahrte Sixel-Bilder auf 16 MiB. Rendering und
Apps benötigen zusätzlichen Speicher. Leere Plätze legen kein Terminal an.
Neun große Apps können den verfügbaren RAM überschreiten.

Temporäre Bilder, Caches und Sitzungen nutzen RAM. Installierte Software,
Dokumente, beim Beenden gespeicherte Einstellungen und App-Dateien nutzen
dauerhaften Speicher. Inkline verhindert keine beliebigen Dateischreibzugriffe.

Das optionale Python-Paket ist unverändertes Python. Für weniger Schreibzugriffe
durch Caches tragen Sie in \code{~/.bashrc} auf dem Tablet Folgendes ein:
\begin{shell}
export PYTHONDONTWRITEBYTECODE=1
export PIP_NO_CACHE_DIR=1
\end{shell}
\code{source ~/.bashrc} ausführen oder eine neue Shell öffnen. Mit
\code{python3 -m pip install --no-compile PACKAGE} vermeiden Sie die zusätzliche
Bytecode-Erzeugung bei pip-Installationen. Pythons isolierter Modus ignoriert
Python-Umgebungsvariablen. Die ausführliche Anleitung steht auf der Website.

Inkline steht unter GPL-3.0-or-later. Pakete enthalten passende Quellarchive
und Hinweise zu Abhängigkeiten. Testberichte unterscheiden automatische
Prüfungen von noch zu bestätigendem Verhalten am Gerät.
